\section{Conditions and limits of delivery}
\label{sec:theory}
We study a known acyclic graph with measured parents. The following statements
explain permissible reuse and possible error amplification. They do not prove
identification for arbitrary neural models or certify the fitted networks.

\begin{proposition}[Mechanism eligibility]
Let $V_i=f_i(V_{\mathrm{Pa}(i)})+U_i$ be an invariant additive-noise
mechanism. Suppose $U_i$ is integrable, independent of the other exogenous
variables, and has mean zero. An intervention policy may depend on previous
experiments, but must choose its action before the current exogenous noises are
realized. Retain a row only when the action does not replace mechanism $i$,
and when retention does not select on its current noise or response conditional
on its parents. Then, in the retained interventional mixture,
\[
 \mathbb{E}[V_i\mid V_{\mathrm{Pa}(i)}=p,\text{retained}]=f_i(p).
\]
Thus squared-loss regression can target this mechanism on supported parent
values, assuming finite second moments. This is an estimand statement; finite
data, optimization and extrapolation guarantees require further assumptions.
\end{proposition}
\begin{proof}
In an acyclic graph, a node's parents depend only on earlier exogenous noises
and the chosen action. The current $U_i$ is independent of those noises and of
past experiments. Since the policy precedes the current realization, and the
retention rule does not condition on $U_i$, the conditional expectation of
$U_i$ remains zero. Substituting the structural equation yields the assertion.
The conditional expectation minimizes population squared loss on the retained
support. No statement outside that support follows.
\end{proof}

Confounding, missing parents, interventions that change rather than clamp a
mechanism, and response-dependent retention can invalidate the premise. In
particular, selecting the best probe using its measured response can change
the noise distribution. All-paid reuse avoids discarding probes on that basis,
but still requires invariant mechanisms and measured parents. The current
confirmation uses a deterministic emulator; this distinction matters for the
prospective noisy systems. A non-additive stochastic SCM also needs a
distributional treatment: composing conditional means need not produce a
downstream interventional mean through a nonlinear mechanism.

\begin{proposition}[Deterministic DAG propagation]
Consider a deterministic SCM and its learned functions $\hat f_i$, evaluated
in topological order under identical correctly clamped interventions.
For each non-clamped node, assume
$|\hat f_i(v_{\mathrm{Pa}(i)})-f_i(v_{\mathrm{Pa}(i)})|\leq\epsilon_i$
at every true parent vector in the deployment support. Assume also
\[
 |\hat f_i(p)-\hat f_i(q)|\leq
 \sum_{j\in\mathrm{Pa}(i)}L_{ij}|p_j-q_j|
\]
on a region containing both true and predicted parent vectors, with
$L_{ij}\geq0$. Writing $e_i=|\hat v_i-v_i|$, one has
\[
 e_i\leq\epsilon_i+\sum_{j\in\mathrm{Pa}(i)}L_{ij}e_j.
\]
Correctly clamped nodes have $e_i=0$. If their rows in $L$ and their local
bounds are set to zero, the vector bound is
$e\leq(I-L)^{-1}\epsilon=\sum_{k=0}^{d-1}L^k\epsilon$,
where $d$ is the number of nodes and $L$ is strictly triangular in topological
order.
\end{proposition}
\begin{proof}
Add and subtract $\hat f_i(v_{\mathrm{Pa}(i)})$ and apply the triangle
inequality. The local approximation bounds one term; the coordinate Lipschitz
inequality bounds the other. Clamps are identical by assumption. Substitution
in topological order gives the path sum. A strictly triangular $d\times d$
matrix satisfies $L^d=0$, so the finite geometric sum is its inverse.
\end{proof}
The bound can amplify a small local error along deep or high-gain paths.
Observed-parent test MSE is neither a uniform $\epsilon_i$ nor a verified
Lipschitz constant. Our residual and support diagnostics are explanatory
measurements, not certified instances of this proposition.

\begin{proposition}[Quantization margin]
For distinct ordered levels $\ell_1<\cdots<\ell_m$, let $Q$ assign the nearest
level, with a fixed tie rule, and let
$b_k=(\ell_k+\ell_{k+1})/2$. Define
$M(y)=\min_k|y-b_k|$. If $|\hat y-y|<M(y)$ then
$Q(\hat y)=Q(y)$. If a pointwise prediction bound $B(y)$ is valid,
$\Pr(Q(\hat y)\ne Q(y))\leq\Pr(M(y)\leq B(y))$.
\end{proposition}
\begin{proof}
The open interval of radius $M(y)$ around $y$ crosses no decision boundary.
Consequently both values occupy the same decision cell. The contrapositive
implies that a quantization error requires $M(y)\leq|\hat y-y|\leq B(y)$.
Strict inequality excludes boundary ties.
\end{proof}
Snapped error depends on the location of errors relative to decision boundaries;
similar MSE alone need not imply similar exact-level accuracy.

\paragraph{Counterexamples.}
For a constant predictor and test response zero, fitting a single training
response zero gives zero risk; adding a noisy response two makes its
least-squares mean one and test risk one. With a noisy training target one,
optimizing a constant from $0.1$ to one reduces training error but increases
test error. For a two-stage chain, suppose the true parent and child are zero
on the observed support, the learned parent predicts $0.1$, and the learned
child is $\hat f(z)=100z$. The child has zero observed-parent error at zero,
but its free-running prediction is ten. These are existence examples, not
descriptions of the measured histories. Analytically solved chains, forks,
colliders and clamps are checked by executable tests.
